\documentclass[11pt,a4paper]{article}

\usepackage[margin=2cm]{geometry}
\usepackage{enumitem}
\usepackage[hidelinks]{hyperref}
\usepackage[T1]{fontenc}
\usepackage[utf8]{inputenc}

\setlength{\parindent}{0pt}
\pagenumbering{gobble}

\begin{document}

%==================== HEADER ====================
\begin{center}
{\Large \textbf{MIGUEL ANGEL ANAY GOMEZ}}\\
Arquitecto de Software | Full Stack Software Developer\\
Lima, Perú\\
Email: \href{mailto:mba.miguel18@gmail.com}{mba.miguel18@gmail.com} \\
Teléfono: +51 992 011 852\\
LinkedIn: \href{https://www.linkedin.com/in/miguel-anay}{linkedin.com/in/miguel-anay}
\end{center}

\vspace{0.4cm}

%==================== PERFIL ====================
\section*{Perfil Profesional}
Ingeniero Industrial con más de 12 años de experiencia en desarrollo de software, arquitectura de aplicaciones y liderazgo técnico. Especialista en diseño de soluciones empresariales, definición de estándares de desarrollo, gestión de cambios y soporte a sistemas críticos. Amplia experiencia coordinando pases a UAT y Producción, asegurando continuidad operativa y cumplimiento de KPIs. Sólido dominio de SQL Server y tecnologías .NET, con experiencia en entornos corporativos, financieros y de seguridad.

%==================== EXPERIENCIA ====================
\section*{Experiencia Profesional}

\textbf{Encora} \\
Software Engineer III \\
Noviembre 2020 – Abril 2025 \\
Cliente: Profuturo AFP
\begin{itemize}[leftmargin=*]
\item Dirección técnica y diseño de soluciones para plataformas web de alta disponibilidad.
\item Definición de arquitectura de aplicaciones alineadas a estándares corporativos.
\item Coordinación de pases a ambientes UAT y Producción.
\item Gestión de cambios, control de versiones y programas fuente.
\item Liderazgo técnico y acompañamiento a equipos de desarrollo.
\item Soporte a aplicaciones críticas durante cierres operativos.
\item Elaboración de documentación técnica, flujos funcionales y diagramas de solución.
\item Tecnologías: React.js, JavaScript, Docker, SQL Server, .NET.
\end{itemize}

\textbf{Prosegur} \\
Analista de Procesos / Desarrollador Web \\
Abril 2018 – Marzo 2020
\begin{itemize}[leftmargin=*]
\item Diseño y validación de soluciones de software para control operativo y gestión de riesgos.
\item Definición de especificaciones técnicas y funcionales.
\item Coordinación con equipos de soporte, operaciones y tecnología.
\item Desarrollo, pruebas y control de calidad de software.
\item Implementación de monitoreo, alertas y soporte a sistemas productivos.
\item Documentación de arquitectura, procesos y modelos UML.
\item Tecnologías: React.js, PHP Laravel, SQL Server, MySQL.
\end{itemize}

\textbf{Prosegur – Seguridad Electrónica} \\
Desarrollador Web / Líder Técnico \\
Abril 2014 – Marzo 2018
\begin{itemize}[leftmargin=*]
\item Diseño de arquitectura de soluciones empresariales para control y monitoreo operativo.
\item Liderazgo de equipos de desarrollo de software.
\item Definición y seguimiento de estándares de desarrollo y bases de datos.
\item Desarrollo en tecnologías .NET y administración de SQL Server.
\item Soporte operativo y mantenimiento de aplicaciones críticas.
\item Tecnologías: C\#.NET, VB.NET, SQL Server, MongoDB, Windows Server.
\end{itemize}

%==================== EDUCACION ====================
\section*{Educación}

\textbf{Universidad Privada del Norte} \\
Diplomado de Gerencia de Tecnologías de la Información con mención en Disrupción Tecnológica \\
Enero 2025 – Septiembre 2025

\vspace{0.2cm}

\textbf{Universidad Nacional de Ingeniería} \\
Ingeniería Industrial \\
2004 – 2012

%==================== ESPECIALIZACION ====================
\section*{Especialización Técnica}

\textbf{Software Developer Full Stack}
\begin{itemize}[leftmargin=*]
\item Arquitectura Hexagonal en C\#.NET – Technical Track (Octubre 2025, 8.5 horas)
\item Visual Basic .NET (Septiembre 2011, 48 horas)
\end{itemize}

\textbf{Data Engineer}
\begin{itemize}[leftmargin=*]
\item Advanced Data Engineer – Diplomado (Septiembre 2025 – Actual, 144 horas)
\item Administración de SQL Server (Junio 2015, 96 horas)
\item Business Intelligence (Febrero 2016, 48 horas)
\end{itemize}

\textbf{Data Science}
\begin{itemize}[leftmargin=*]
\item Python for Data Science (Octubre 2025, 20 horas)
\item Gen AI Training Path – Technical Track (Marzo 2025, 60 horas)
\end{itemize}

\textbf{DevOps Engineer}
\begin{itemize}[leftmargin=*]
\item Azure Fundamentals AZ-900 (Noviembre 2024, 9 horas)
\end{itemize}

%==================== COMPETENCIAS ====================
\section*{Competencias Técnicas}

Arquitectura de Software, Gestión de Cambios, Liderazgo Técnico, UML, Documentación Técnica, KPIs, Soporte a Producción, SQL Server, .NET, Control de Versiones.

%==================== HERRAMIENTAS ====================
\section*{Herramientas}

SQL Server, .NET, React.js, Docker, Git, Microsoft Project, UML, Visio, Azure, Windows Server.

\end{document}